\documentclass[11pt,a4paper]{article}

\usepackage[margin=1in]{geometry}
\usepackage{xcolor}
\usepackage{listings}
\usepackage{amsmath}
\usepackage{enumitem}
\usepackage[hidelinks]{hyperref}

\definecolor{codebg}{gray}{0.96}
\definecolor{decided}{RGB}{0,110,60}
\definecolor{open}{RGB}{170,60,0}

\lstdefinestyle{json}{%
  basicstyle=\ttfamily\small,
  backgroundcolor=\color{codebg},
  columns=fullflexible,
  breaklines=true,
  showstringspaces=false,
  frame=none,
  xleftmargin=1em,
  framexleftmargin=1em,
  aboveskip=0.6em, belowskip=0.6em
}

\lstdefinestyle{ml}{%
  language=[Objective]Caml,
  basicstyle=\ttfamily\small,
  backgroundcolor=\color{codebg},
  keywordstyle=\color{blue!60!black},
  commentstyle=\itshape\color{black!45},
  columns=fullflexible,
  breaklines=true,
  showstringspaces=false,
  frame=none,
  xleftmargin=1em,
  framexleftmargin=1em,
  aboveskip=0.6em, belowskip=0.6em
}

\providecommand{\llbracket}{[\![}
\providecommand{\rrbracket}{]\!]}

\newcommand{\pending}{\textcolor{open}{\textsf{[pending]}}}
\newcommand{\yes}{\textcolor{decided}{\textsf{decided}}}
\newcommand{\selected}{\textcolor{decided}{$\leftarrow$ \textbf{selected}}}
\newcommand{\decision}[1]{%
  \par\medskip\noindent\textcolor{decided}{\textbf{Decision.}}\quad #1\par\medskip}
\newcommand{\rationale}[1]{%
  \par\noindent\textbf{Rationale.}\quad #1\par\medskip}
\newcommand{\limitation}[1]{%
  \par\noindent\textbf{Known limitation.}\quad #1\par\medskip}
\newcommand{\implnote}[1]{%
  \par\noindent\textbf{Implementation note.}\quad #1\par\medskip}
\newcommand{\options}{\paragraph{Options considered.}}

\setlist[itemize]{itemsep=0.15em, topsep=0.35em}
\setlist[enumerate]{itemsep=0.5em, topsep=0.4em}

\newcommand{\proposed}{\textcolor{open}{\textsf{proposed}}}
\newcommand{\proposal}[1]{%
  \par\medskip\noindent\textcolor{open}{\textbf{Proposal.}}\quad #1\par\medskip}
\newcommand{\consequence}[1]{%
  \par\noindent\textbf{Consequence.}\quad #1\par\medskip}

\title{\textbf{Tatami --- Phase 3}\\[0.3em]
       \large Composable Kernels: What Has To Line Up}
\author{Durwasa Chakraborty}
\date{Design document --- first draft, nothing decided}

\begin{document}
\maketitle

\section{Purpose and scope}

Phase~3 takes its architecture from BOSS: a query engine built out of separable
kernels, each free to run on a different device, with a plan split across them
at run time rather than compiled for one target. The design note acknowledges
that debt and does not restate the architecture, and neither does this
document.

What this document does is different, and is the reason it exists. Adopting a
kernel decomposition is only half a decision. The other half is the set of
things that must be true of \emph{our} generated code for that decomposition to
be applicable at all --- what a kernel is when the operators are effect
handlers rather than library functions, what crosses the boundary between two
of them, what makes one legal on a device, and which of the twenty-two
decisions already taken in Phase~1 happen to supply those conditions and which
happen to obstruct them.

Every proposal below is offered to be argued with. Phase~1 earned the right to
write \textsf{decided} next to its scenarios by exhausting the alternatives
first; nothing here has been through that, and the status column says so.

\paragraph{The unit is the fragment.} Throughout, a \emph{fragment} is a
region of a query's lowering that could in principle be dispatched somewhere.
The whole question is what delimits one and what may be said about it, so the
word is deliberately loose until \S3 makes it precise.

\section{Status}

\begin{center}
\begin{tabular}{@{}llp{4.6cm}@{}}
\hline
\textbf{Ref} & \textbf{Question} & \textbf{Status} \\
\hline
K1 & What a kernel is                     & \proposed \\
K2 & What crosses the boundary            & \proposed \\
K3 & The effect signature                 & \proposed \\
K4 & The dispatch predicate               & \proposed \\
K5 & Composition and transfer placement   & \proposed \\
K6 & What costing is still for            & \proposed \\
K7 & The invariant                        & \proposed \\
\hline
\end{tabular}
\end{center}

\noindent
Sections \S10 and \S11 are not questions. They record what Phase~1 already
supplies and what it obstructs, and they are the part of this document least
likely to change, because they are readings of a settled specification rather
than proposals.

\section{K1. What a kernel is}

BOSS's kernels are libraries --- Velox, ArrayFire --- and a plan is split
across them by partial evaluation. We have no libraries. Phase~2 leaves us
something else: generated accessors that \texttt{perform} rather than return,
and a hole where the handler goes. The design note is explicit that what the
handlers do is ``deliberately left open here''. This is that question.

\options
\begin{enumerate}[label=\textbf{(\alph*)}]
  \item \textbf{The kernel is a handler.} A kernel is a handler for some set
        of effect constructors, together with the fragment running under it.
        Dispatching a fragment to a different device means installing a
        different handler over the same code. \selected
  \item \textbf{The kernel is a physical operator}, in the ordinary database
        sense --- a hash join, a scan --- and effects are merely how operators
        request rows from each other.
  \item \textbf{The kernel is a generator pipeline stage}, delimited by where
        one \texttt{next} feeds another.
  \item \textbf{The kernel is a whole relation's traversal}, so the unit of
        dispatch is a table.
\end{enumerate}

\proposal{A kernel is a handler plus the fragment it handles. Device
selection is handler selection: the fragment's own code is unchanged, and what
changes is the interpretation installed over it.}

This is the only option that costs nothing to express, because Phase~2 has
already built it. The design note's §3.2 claims that a polling fragment can be
placed ``on another core, another node, or a GPU without changing the code that
consumes it'', and under (a) that claim is not an aspiration but a restatement
of what an effect handler is. Two handlers for \texttt{Yield\_qty} --- one
walking a CPU-resident column, one enqueueing a GPU kernel and draining its
output --- are interchangeable at the point of installation, and the consumer
cannot tell which it got.

\rationale{Options (b) and (c) are not wrong so much as premature: they name a
unit that we would then have to \emph{relate} to the handler structure Phase~2
generates, and that relation is extra work with no evident payoff. Option (d)
is too coarse. A single relation can contain both a fixed-width scalar column
that a GPU handles well and a variable-width string column it handles badly,
and dispatching the table as one unit forces the worse device on the better
column.}

\limitation{Handler boundaries are a property of how the query was written,
not of what would be good to dispatch. A consumer that installs one handler
around its entire computation presents exactly one kernel, and no
decomposition is available at all. This makes the granularity of dispatch
depend on the query language of K1's sibling question --- the language must
generate handler boundaries where fragments should be, rather than leaving
them where a person happened to put them.}

\section{K2. What crosses the boundary}

BOSS composes its kernels over ``a unified exchange format for plans and
data''. Kernels that cannot agree on a representation cannot compose, so this
is prior to everything about dispatch: it is what makes two kernels the same
shape of thing.

\proposal{The exchange format is a \emph{column bundle}: for each relation, one
dense typed array per column; for each nullable column, a validity mask
alongside it; and for a child relation, its \texttt{parent\_id} and
\texttt{idx} columns exposed as ordinary arrays like any other.}

Nothing in that sentence is invented for Phase~3. Every clause is a reading of
a Phase~1 decision, which is the substance of \S10.

\consequence{A kernel's input and output are both column bundles, so the
composition of two kernels is defined whenever the second's expected columns
are a subset of the first's produced ones. No format negotiation, no tagging,
and no per-element dispatch.}

\limitation{A column bundle is a materialized representation, and the design
note's §3.2 argues against materialization --- a generator's whole virtue is
that its live state is one element and an accumulator. The two are reconciled
only if a bundle is understood as the unit of \emph{transfer between kernels}
and never as the unit of consumption within one: a kernel streams internally
and materializes at its edges. Whether that is achievable, or whether GPU
residency forces materialization deeper than the edges, is not settled here
and is the most likely place for this proposal to fail.}

\section{K3. The effect signature}

The design note asks whether dispatch can be derived from ``the effect
signature of the generated accessors''. Before that can be answered, the
signature has to exist, and one fact governs how:

\paragraph{OCaml 5 does not track effects in the type system.} The PLDI 2021
retrofit shipped handlers without effect types. \texttt{Effect.perform} is
invisible in a function's type, and nothing about which effects a fragment
performs can be read off a signature the compiler checked. The project's
subtitle, ``effect-typed OCaml'', names something we construct, not something
we are given.

\options
\begin{enumerate}[label=\textbf{(\alph*)}]
  \item \textbf{Infer it} --- an effect analysis over the generated code,
        recovering for each fragment the set of constructors it may perform
        without handling.
  \item \textbf{Generate it} --- emit the signature alongside the module,
        since Phase~1 generated the accessors from the schema and therefore
        already knows what each one performs. \selected
  \item \textbf{Declare it} --- the query language requires the author to
        annotate, and the annotation is checked.
\end{enumerate}

\proposal{Generate it. The signature is an output of the same pass that emits
the accessor, derived from the schema rather than recovered from the code.}

\rationale{This is what the design note means by ``the provenance of every
relation is known: for each file we can determine which code is responsible for
populating which part of the database''. Provenance is available because we
generated the code; the effect signature is the same fact in a different
notation, and recovering by analysis what we could have emitted directly is
work performed against ourselves.}

\paragraph{The vocabulary.} A first cut, stated concretely so that it can be
attacked. Phase~2 emits one constructor per field; classify each by what its
payload costs a device, not by which field it came from:

\begin{center}
\begin{tabular}{@{}llp{6.2cm}@{}}
\hline
\textbf{Class} & \textbf{Arises from} & \textbf{What it says} \\
\hline
$E_{\mathrm{fix}}(\tau)$ & \texttt{int}, \texttt{float}, \texttt{bool}
  & Fixed-width scalar. One array, one stride. \\
$E_{\mathrm{var}}$ & \texttt{string}
  & Variable width. Needs offsets plus bytes. \\
$E_{\mathrm{opt}}(E)$ & any \texttt{option} column
  & As $E$, plus a validity mask. \\
$E_{\mathrm{ref}}(M)$ & a foreign key
  & A dependent read into $M$: a gather, not a scan. \\
$E_{\mathrm{rec}}$ & S19's self-referencing key
  & A dependent read whose depth is not statically bounded. \\
\hline
\end{tabular}
\end{center}

\noindent
A fragment's signature is the multiset of classes it may perform without
handling. Note that the classification is by \emph{layout consequence}: two
fields of different names and different meanings but the same width are the
same class, because a device cannot tell them apart either.

\limitation{Generated signatures cover generated code and nothing else. A
query that computes anything of its own --- an arithmetic predicate, a
user-written fold --- performs no generated effect, so its signature is empty
and the predicate of K4 says nothing useful about it. This is the sharpest
constraint on the query language: either it is built only from combinators
whose signatures we also generate, or K3 must fall back to option (a) for the
parts it does not cover.}

\section{K4. The dispatch predicate}

With K3's vocabulary the question of §4.1 --- which target a fragment ``can
legally be lowered to'' --- becomes a predicate on a set rather than an
estimate over statistics.

\proposal{A fragment is GPU-legal exactly when its signature contains no
$E_{\mathrm{rec}}$, every $E_{\mathrm{ref}}(M)$ it contains targets a relation
already resident on the device, and every $E_{\mathrm{var}}$ it contains is
handled by a kernel that accepts the offsets-plus-bytes encoding. Otherwise it
is CPU-only.}

Read clause by clause, each says something a device architect would recognise:

\begin{itemize}
  \item \textbf{No $E_{\mathrm{rec}}$.} Recursion is unbounded pointer
        chasing with data-dependent depth. It is the one class that cannot be
        given a static iteration count, and it is where GPUs are worst.
  \item \textbf{Resident $E_{\mathrm{ref}}$.} A foreign key is a gather. A
        gather into device memory is a normal GPU operation; a gather into
        host memory across the bus, per element, is not. So the clause is not
        ``no joins'' but ``no joins whose target is elsewhere'', which is a
        placement question and therefore composes with K5.
  \item \textbf{Encoded $E_{\mathrm{var}}$.} Strings are not excluded, only
        made conditional on the kernel that handles them. A dictionary-encoded
        or offsets-encoded string column is workable; a column of independently
        allocated OCaml strings is not.
\end{itemize}

\rationale{The predicate is decidable on the signature alone, takes no
statistics as input, and returns the same answer for a corpus of ten documents
and ten million. That is what distinguishes it from a cost model and is the
whole content of the design note's question.}

\limitation{Decidable is not the same as useful. The predicate says what is
\emph{legal}, and legality is a weak claim: a fragment can be GPU-legal and
still run worse there, because the predicate has no notion of how many rows
will flow through it. That is the honest reading of §4.1's own word,
``legally'', and it means the strong form of the design note's question ---
placement \emph{instead of} costing --- is not established by this predicate.
See K6.}

\section{K5. Composition and transfer placement}

Two kernels compose when the second handles what the first leaves unhandled
and their column bundles agree. If they are placed on different devices, data
has to move, and where it moves is usually the dominant cost in a
heterogeneous plan.

\proposal{A transfer is inserted exactly at a handler boundary where the
device changes. There is no other place one can occur.}

This is the structural payoff of K1. Because a kernel \emph{is} a handler, and
because a device is chosen when a handler is installed, a device change and a
handler change are the same event. Transfer points are therefore not
discovered by analysis and not estimated: they are enumerable by walking the
handler structure of the plan.

\consequence{Two useful things follow. The number of transfers in a plan is a
syntactic property of that plan, available before anything is executed or
costed. And minimising transfers becomes a question of \emph{where to put
handler boundaries}, which is a rewriting problem over a known structure rather
than a search over an unknown one.}

\limitation{This makes the plan space smaller than it should be. A transfer
that would be profitable in the middle of a long fragment cannot be expressed,
because there is no handler boundary there to attach it to, and introducing one
means splitting the fragment --- which changes what K4 says about each half.
Splitting is therefore a plan transformation in its own right, and the
interaction between splitting and the predicate is not worked out here.}

\section{K6. What costing is still for}

§4.2 of the design note says plans are enumerated and costed; §6 asks whether
dispatch can be derived from effect signatures \emph{rather than} from a cost
model. Taken together these are in tension, and the tension has to be resolved
before §5's differential evaluation measures anything interpretable. Three
readings, in increasing ambition:

\begin{enumerate}[label=\textbf{(\arabic*)}]
  \item \textbf{Legality only.} The signature filters the candidate targets;
        the cost model chooses among survivors. \selected
  \item \textbf{Placement, costing for the rest.} The signature fixes the
        device outright; costing only picks join implementations and physical
        operators once it has.
  \item \textbf{Placement instead of costing.} No cost model participates in
        placement at all.
\end{enumerate}

\proposal{Claim (1), and say so plainly. K4's predicate decides which targets
a fragment may be lowered to; a cost model over the operator repository chooses
among them. The contribution is that the legality half is derived
structurally, and does not need statistics.}

\rationale{(1) is what §4.1 literally says, and it is the only one of the
three that K4 as stated actually supports, since the predicate is blind to
cardinality. Claiming (3) and then quietly consulting a cost model would make
§5's comparison unfalsifiable --- there would be no configuration of the system
that could fail to confirm the claim. Claiming (1) and demonstrating it costs
less and is worth more.}

\limitation{Claim (1) is a weaker result than §6's phrasing suggests, and it
should be conceded in the write-up rather than blurred. What remains genuinely
ours under (1) is still a real claim: that the legality half of dispatch is a
compile-time property of generated code, and that a system deriving it this way
needs no statistics to know what it \emph{may} do.}

\paragraph{Where the statistics would come from, if (2) is ever attempted.}
S20 of Phase~1 already computes per-path observation counts --- present with a
value, explicitly null, absent --- and declares them diagnostic only. They are
selectivity statistics under another name. Promoting them from diagnostic to
input would make placement a function of the corpus, which K7 is precisely an
attempt to avoid, so this is recorded as an available move and not taken.

\section{K7. The invariant}

Phase~1 holds together because almost every decision in it defends one
property: \emph{the schema is a function of document structure, not of data
values}. It is what rules out heuristic map detection, thresholds on
optionality, recursion detection, and collision-suffix naming. A phase without
such a property is a pile of local choices.

\proposal{The Phase~3 counterpart: \emph{legality is a function of the effect
signature, not of the data}. A fragment's set of admissible targets does not
change because the corpus grew.}

The two compose into a single statement spanning the project:

\begin{center}
\emph{The schema is a function of the documents' structure, and the plan space
is a function of the schema.}
\end{center}

\noindent
Neither half mentions data values, so neither does the composition. A corpus
that doubles changes no schema, and therefore changes no fragment's legal
targets --- only which of them is chosen.

\consequence{This is testable, and cheaply. Run inference and lowering over a
corpus, then over a corpus ten times the size with the same structure, and
compare the plan spaces. They must be identical. It is a weaker property than
§5's differential correctness oracle and far cheaper to check, which makes it
a good early regression test rather than an end-of-project one.}

\limitation{The invariant holds only for the legality half, by construction:
under K6's claim (1) the \emph{chosen} plan does depend on statistics and
therefore on data. So the honest statement is that the plan \emph{space} is
data-independent while the selection within it is not --- which is a sharper
and more defensible claim than one about placement as a whole, and is the form
that should appear in the write-up.}

\section{What Phase 1 already supplies}

The alignment worth checking is not with BOSS but with our own front end. A
kernel decomposition needs a data representation with particular properties,
and several Phase~1 decisions --- taken for reasons that had nothing to do with
devices --- turn out to supply exactly those.

\begin{center}
\begin{tabular}{@{}lp{5.0cm}p{5.4cm}@{}}
\hline
\textbf{Ref} & \textbf{Decided, for its own reasons} & \textbf{What it gives K2} \\
\hline
S16, S18 & Reject a path whose observed types have no least upper bound.
  & Every column has exactly one type. A dense typed array with no per-element
    tag and no branch in the inner loop. \\
S1a & Widen \texttt{int} to \texttt{float} on mix.
  & Every column has one width. Uniform stride. \\
S4, S16b & Nullability is a flag on the column, not a point in the type
    lattice.
  & Validity is separable from values: a mask beside a dense array, which is
    the layout a GPU wants. \\
S13 & An ordering index is required, not optional.
  & Array order survives as an ordinary integer column, so a sequence is
    reconstructible from an unordered kernel result. \\
S14 & A scalar array is a table with no module.
  & The GPU-shaped case arrives with nothing wrapped around it: parent key,
    index, one value column. \\
S21.6 & \texttt{parent\_id} and \texttt{idx} are public record fields.
  & Join keys are ordinary columns of the bundle, not hidden state a kernel
    would have to be told about. \\
S22 & Table names are path-qualified and injective.
  & Every relation has a stable name to be addressed by in a plan. \\
\hline
\end{tabular}
\end{center}

\paragraph{The one worth dwelling on.} S16 and S18 reject heterogeneous paths,
and the Phase~1 argument for doing so is entirely about data-dependence: a
generated variant would gain a constructor when a document was added. The
device consequence was not part of that argument and is nonetheless large. A
union column requires a per-element tag and a branch on it, which is close to
the worst thing one can put in a GPU kernel; rejecting unions means no column
ever needs one. A decision justified on compiler grounds turns out to be a
precondition for the systems half, which is the strongest evidence available
that the two phases are pulling the same direction.

\section{What Phase 1 obstructs}

Recorded plainly, because the parts that do not line up are more informative
than the parts that do. Several are flagged in \texttt{lowering.tex} itself as
Phase~3's problem.

\begin{itemize}
  \item \textbf{S21.5 --- \texttt{val all : db -> t list}.} Materializes an
        entire table, which is what §3.2 of the design note argues against and
        what K2 permits only at kernel edges. Phase~1 names \texttt{t Seq.t}
        as the revision and notes it breaks consumers rather than extending
        them. Under K1 there is a third option Phase~1 could not see: the
        entry point performs an effect like every other accessor, and the
        handler decides whether to materialize.
  \item \textbf{S21.2 --- the explicit \texttt{db} handle.} Its own known
        limitation names option (d), ``the store is the effect handler,
        supplying rows in response to a performed effect'', as the eventual
        destination. K1 makes that destination the whole design rather than an
        eventual refinement, so this is the Phase~1 decision Phase~3 most
        wants revisited.
  \item \textbf{S19 --- recursion.} $E_{\mathrm{rec}}$ is CPU-only under K4,
        so a corpus that is mostly trees --- comment threads, file systems,
        organisation charts --- gets no GPU benefit at all. This is a real
        limit on which workloads can demonstrate anything, and it should
        inform which corpora §5 evaluates over.
  \item \textbf{S1b --- out-of-range integers widen to \texttt{float}.}
        Precision is already lost at ingest, before any kernel runs. If GPU
        kernels reduce in single precision, that is a second and independent
        loss on the same column, and §5's correctness oracle would see the two
        plans disagree for reasons that have nothing to do with placement.
        Comparison over such columns needs a tolerance, and saying so in
        advance is better than discovering it as a failing test.
  \item \textbf{S7 --- every nested field costs a join.} Phase~1's own known
        limitation observes that flattening ``matters for Phase~3 where a
        flattened column can be reached without one''. Under K3 the difference
        is exactly $E_{\mathrm{fix}}$ against $E_{\mathrm{ref}}$: flattening a
        total nested object turns a gather into a scan, and therefore turns a
        residency condition into no condition at all. This is the one place
        Phase~3 has a concrete reason to ask Phase~1 to reopen a decision.
\end{itemize}

\section{Not addressed}

Named so that their absence is deliberate rather than overlooked.

\begin{itemize}
  \item \textbf{The query language.} K3's limitation constrains it severely
        --- signatures cover generated code only --- but nothing here chooses
        between a combinator library, a PPX quotation, and analysis of
        ordinary consumer code.
  \item \textbf{Storage.} K2 describes what crosses a kernel boundary and says
        nothing about what sits behind \texttt{db}. Phase~1 declined the
        question as being about storage rather than lowering; Phase~3 cannot,
        since a row-major store cannot produce a column bundle without a
        transposition that would dominate everything measured.
  \item \textbf{Which kernels.} Whether we compose existing ones (ArrayFire,
        Velox, as BOSS does) or emit our own (OpenCL, as Voodoo does). K4's
        predicate is stated over layout consequences and is neutral between
        them, which is deliberate.
  \item \textbf{The baseline.} §5 compares against ``the plan as written'',
        and that phrase needs a definition sharp enough to implement and
        honest enough to be fair. A baseline chosen to lose is not evidence.
\end{itemize}

\end{document}
